\chapter{Literature Review}\label{ch:A}

\section{Photoelectrochemical Water Splitting: Fundamentals}

\subsection{Semiconductor Photoelectrodes}

The efficiency of that process is determined by several competing elements ---
the amount of light absorbed by the material, the time it takes for minority
charge carriers to reach the surface before they recombine, and the ability
of the surface to catalyse the water oxidation reaction~\cite{gratzel2001photoelectrochemical}.

For unassisted water splitting (i.e., no external voltage applied) there are
a number of simultaneous requirements that the material must
meet~\cite{bolton1978limiting}: the band gap should fall in approximately the
range 1.8--2.2\,eV --- broad enough to generate sufficient thermodynamic driving
forces for both half-reactions, but narrow enough to allow for absorption of a
useable fraction of the solar spectrum.
The band-edge positions must be positioned such as to bracket the water oxidation
and reduction potentials.
Finally, the material must remain chemically stable in the electrolyte under
prolonged irradiation.
Thus far, there exists no single known material that meets all of the above
criteria simultaneously; hence, why the experimentally accessible phase-space
continues to be so large~\cite{walter2010solar}.

\subsection{Key Performance Metrics}

Three metrics are universally reported in PEC literature and constitute the
prediction targets in this work.

\paragraph{Photocurrent Density $J_{ph}$} (mA/cm\textsuperscript{2}) is measured
by linear sweep voltammetry under AM~1.5G illumination at a standardised intensity
of 100\,mW/cm\textsuperscript{2}.
It reflects the rate of charge transfer to the electrolyte and serves as the
primary proxy for combined solar absorption and charge collection efficiency.
State-of-the-art values exceed 6\,mA/cm\textsuperscript{2} for
BiVO\textsubscript{4}~\cite{kim2014bismuth} and reach above
20\,mA/cm\textsuperscript{2} for GaN/InGaN nanowire photoanodes~\cite{kibria2015multibandgap}.

\paragraph{Solar-to-Hydrogen Efficiency (STH)} is the gold standard for overall
device efficiency, defined by Equation~\eqref{eq:sth}.
The Shockley--Queisser-analogous theoretical maximum for a single-junction device
with $E_g \approx 2.0$\,eV is approximately 16.8\%~\cite{walter2010solar}.
The highest experimentally verified STH values for single-junction photoelectrodes
under unassisted, unbiased conditions are approximately 12\% for III-nitride tandem
structures~\cite{kibria2015multibandgap}.

\paragraph{Hydrogen Evolution Rate (HER)} ($\mu$mol\,h\textsuperscript{-1}\,cm\textsuperscript{-2})
is determined by gas chromatography.
It is related to photocurrent through Faraday's law:
\begin{equation}
  \mathrm{HER} = \frac{J_{ph} \times 3600}{2 F} \times \eta_F \times 10^{6}
  \label{eq:her}
\end{equation}
where $F = 96\,485$\,C/mol is Faraday's constant.
HER is reported less frequently than $J_{ph}$ because it requires a gas-tight
electrochemical cell and calibrated gas chromatograph, which explains the severe
sparsity of this target in publicly available data.

\subsection{Principal Material Systems}

Table~\ref{tab:materials} summarises the photoelectrode materials covered in this
study, their typical bandgap ranges, and representative photocurrent densities.

\begin{table}[htb]
\centering
\caption{Common PEC photoelectrode materials, bandgap ranges, and typical
         photocurrent performance.}
\label{tab:materials}
\begin{tabular}{p{3.8cm}ccc}
\toprule
Material & Bandgap (eV) & Typical $J_{ph}$ (mA/cm\textsuperscript{2}) & Role \\
\midrule
BiVO\textsubscript{4}       & 2.4--2.5 & 1--7   & Photoanode  \\
$\alpha$-Fe\textsubscript{2}O\textsubscript{3} & 2.0--2.2 & 0.3--4 & Photoanode \\
WO\textsubscript{3}         & 2.6--2.8 & 1--6   & Photoanode  \\
TiO\textsubscript{2}        & 3.0--3.2 & 0.2--2 & Photoanode  \\
ZnO                         & 3.2--3.4 & 0.5--3 & Photoanode  \\
Cu\textsubscript{2}O        & 1.9--2.1 & 1--5   & Photocathode \\
GaN / InGaN                 & 0.7--3.4 & 5--22  & Photoanode  \\
Ta\textsubscript{3}N\textsubscript{5} & 2.0--2.1 & 0.3--4 & Photoanode \\
Si                          & 1.1      & 8--30  & Photocathode \\
\bottomrule
\end{tabular}
\end{table}

\section{Machine Learning for Materials Discovery}

\subsection{General Approaches}

Materials informatics applies statistical and machine learning methods to accelerate
the discovery of new functional materials~\cite{ramprasad2017machine}.
Pilania et al.~\cite{pilania2013accelerating} demonstrated that random forest models
trained on DFT-computed features could predict perovskite formation energies with
mean absolute errors below 0.1\,eV/atom --- a task that required only minutes of
computation compared to hours of DFT.
Isayev et al.~\cite{isayev2017universal} introduced universal fragment descriptors
capable of predicting band structures across thousands of inorganic compounds.

In the photocatalysis domain, Schmidt et al.~\cite{schmidt2019recent} reviewed
supervised regression models for bandgap prediction from composition alone, reporting
$R^2$ values of 0.85--0.94 on large DFT-computed datasets of thousands of entries.
However, such accuracy relies on uniform, high-quality computed labels; experimental
datasets with heterogeneous measurement conditions achieve notably lower $R^2$
values even with sophisticated models~\cite{goldsmith2018machine}.

\subsection{ML Applied to PEC Systems}

Compared with solid-state property prediction, PEC performance prediction is
inherently more challenging because the output depends not only on the semiconductor
material but also on electrolyte chemistry, electrode geometry, surface catalysts,
and illumination conditions --- factors that vary significantly across different
research groups.

Wu et al.~\cite{wu2019machine} trained gradient-boosted models on 150 published
oxygen evolution catalyst measurements, demonstrating that material composition alone
could predict overpotential with a mean absolute error of 27\,mV.
For photocatalytic water splitting, Chen et al.~\cite{chen2019graph} applied graph
neural networks to predict the hydrogen evolution rate of organic photocatalysts
from molecular structure, achieving $R^2 = 0.61$ on a test set of 100 compounds.
Neither of these studies attempted cross-material generalisation --- each was built
around a single material family by design.
A model trained exclusively on transition-metal oxide catalysts carries no information
about organic photocatalysts, and vice versa.
Our dataset spans nine material systems deliberately, which makes the modelling harder
but produces predictions that remain useful even when the researcher has not yet
decided which material class to focus on.

\subsection{Ensemble Tree Methods and Model Justification}

\subsubsection{Why Tree-Based Ensembles?}

The dataset size made the model choice fairly straightforward.
With 400--560 training rows per target, neural networks are not a realistic option.
A modest two-layer MLP with 128 units per layer has over 16,000 parameters ---
that is roughly 30 parameters per training sample for photocurrent,
a ratio at which overfitting is essentially guaranteed regardless of regularisation.
We verified this in preliminary experiments: MLP validation loss diverged within
a few epochs regardless of dropout or weight decay settings.
Tree-based ensembles do not have this problem.
A 2022 benchmark by Grinsztajn et al.~\cite{grinsztajn2022tree} across 45 datasets
confirmed that gradient-boosted trees consistently match or outperform deep learning
on tabular data at this scale.

There are also practical advantages specific to this dataset.
Decision trees handle missing values natively, which matters when 51\% of bandgap
values and 79\% of temperature values are absent in the original data.
They require no feature scaling.
Their importance scores also give a direct way to check whether the
physics-motivated features we engineered are actually being used ---
which turned out to be an important validation step.

\begin{itemize}
  \item \textbf{Robustness to missing values and mixed types.}
        Decision trees natively partition on any threshold, making them
        insensitive to feature scaling and partially missing inputs --- both
        common in experimental datasets.
  \item \textbf{Low sample complexity.}
        A random forest with 300 trees can generalise well with as few as
        200 training samples; deep networks typically require orders of
        magnitude more data to avoid overfitting.
  \item \textbf{Built-in feature selection.}
        Splitting criteria automatically emphasise informative features and
        ignore irrelevant ones, providing interpretable importance scores.
  \item \textbf{Bias--variance behaviour.}
        Bagging (Random Forest) reduces variance at the cost of a slight
        increase in bias; boosting (XGBoost) progressively reduces bias by
        fitting residuals. Their combination in an ensemble provides
        complementary error reduction.
\end{itemize}

\subsubsection{XGBoost}

XGBoost~\cite{chen2016xgboost} builds an additive ensemble of decision trees by
sequentially fitting residuals with a regularised objective:
\begin{equation}
  \mathcal{L}(\phi) =
    \sum_{i} \ell\!\left(y_i,\, \hat{y}_i\right)
    + \sum_{k} \Omega(f_k)
  \label{eq:xgb_loss}
\end{equation}
where $\ell$ is a differentiable loss function, $\hat{y}_i$ is the current
ensemble prediction, and the regularisation term
$\Omega(f_k) = \gamma T + \tfrac{1}{2}\lambda \|w\|^2$
penalises both the number of leaves $T$ and the leaf weight magnitude.
This explicit regularisation reduces overfitting on small datasets compared to
classical gradient boosting~\cite{friedman2001greedy}.

\subsubsection{Random Forest}

Random Forest~\cite{breiman2001random} averages predictions from $B$
independently trained decision trees, each fitted on a bootstrap sample of the
training data with a random subset of features considered at each node split.
The resulting de-correlation of individual trees substantially reduces the
variance of the ensemble compared to a single deep tree.
Beyond point prediction, the spread across individual trees provides a natural
uncertainty proxy~\cite{wager2014confidence} that is exploited in the confidence
scoring scheme of this work (Section~\ref{ch:B}).

\subsubsection{Alternatives Considered}

\paragraph{Linear models} (Ridge, LASSO) were explored and discarded.
We first tried using Ridge regression as a ``sanity'' check.
The results showed that the Ridge model achieved an $R^2$ value of 0.11 on
photocurrent --- only slightly above simply predicting the mean of the training set.
This result did not come as a surprise: the relationship between bandgap energy
and photocurrent is very non-linear.
For example, a 0.1\,eV shift in bandgap from the optimum of 2.0\,eV causes large
changes in the amount of light absorbed by the material, while the same 0.1\,eV
shift near 3.2\,eV has little or no effect on light absorption.
There are too many combinations of polynomial and interaction terms to add manually
to make Ridge regression usable without creating an overly complicated model.

\paragraph{Neural networks} were eliminated because we did not have sufficient data.
Even when applying both aggressive dropout and weight decay, a simple two-layer
multi-layer perceptron (MLP) would begin to fit noise in the data after just a few
epochs on our training sets.
Using 400--560 training samples with a relatively small MLP of two hidden layers
each containing 128 units created a total of 16,000+ weights --- a ratio of
approximately 30:1 of weights to training samples.
This high ratio was confirmed through empirical testing: regardless of the strength
of regularisation applied, validation loss began to diverge from training loss
during the first several epochs.

\paragraph{Support Vector Regression (SVR)} with an RBF kernel was able to perform
better than Ridge regression but lost 0.08--0.12 $R^2$ units compared to XGBoost
on each target variable.
In addition to requiring standardisation of all input variables, SVR also requires
that the standardiser be saved as part of the prediction pipeline --- failure to do
so leads to silent failures at inference time.
Since XGBoost outperformed SVR consistently and without needing standardisation,
there was no reason to continue using SVR in the final prediction pipeline.

\section{Feature Engineering for PEC}

Input data is limited.
The information contained within individual values of bandgap and pH does contain
some level of value; however, the Nernst equation describes how the two factors
combine to create the actual chemical potential --- the flat-band potential of a
metal oxide shifts with a drop of $-59$\,mV per unit of pH.
Therefore, a material with a large bandgap will behave significantly differently
in an acid versus a base environment.
If the model can only view bandgap and pH as independent features then it would
need to find out about their interaction through the first few hundred training
examples.
A direct shortcut into this relationship has been engineered using the product of
these two as a new feature.
Huo et al.~\cite{huo2019unified} demonstrated a 30\% improvement in predicting
photocatalytic activity from using a similar method.

For PEC systems, the key electrochemical relationships governing performance are:

\begin{enumerate}
  \item The \textbf{Nernst equation}: the flat-band potential of a metal oxide
        semiconductor shifts at $-59$\,mV per pH unit, so the net thermodynamic
        driving force depends on the product $E_g \times \text{pH}$.
  \item The \textbf{Butler--Volmer relation}: the reaction rate scales
        exponentially with overpotential $\eta = V_{bias} - (E_g - 1.23)$,
        motivating the feature \texttt{overpotential\_proxy}.
  \item \textbf{Faraday's law}: the H\textsubscript{2} evolution rate is
        directly proportional to photocurrent, providing a deterministic
        relationship exploited during data augmentation
        (Equation~\eqref{eq:her}).
\end{enumerate}

These physical relationships justify the interaction features described in
Chapter~\ref{ch:B}.

\section{Gap Analysis}

A review of the existing literature reveals the following gaps that motivate
the present work:

\begin{enumerate}
  \item No publicly available, ML-ready dataset covers a broad range of
        inorganic PEC photoelectrode materials with experimental performance
        metrics for all three targets (Photocurrent, STH, HER) simultaneously.

  \item Existing PEC ML studies do not address the severe data sparsity of
        STH and HER targets --- the two most practically relevant metrics
        for assessing solar-to-fuel efficiency.

  \item Deployed, interactive web applications that expose PEC prediction
        models to experimental researchers have not been reported in the
        literature.

  \item Uncertainty quantification in PEC property prediction --- including
        model confidence scores and inter-model agreement --- is absent from
        prior work, limiting the practical usefulness of existing predictions.
\end{enumerate}

This work addresses all four gaps through a combination of data engineering,
model design, and system deployment.
